%% 自動生成（tools/mkchecklist.py）。各章の「ふりかえり」の項目を集めたもの
\begingroup\footnotesize
\def\bkcb{\tikz[baseline=-0.5ex]\draw[\bkTF, line width=0.7pt, rounded corners=0.5pt] (0,-0.3em) rectangle (0.8em,0.5em);}
\begin{longtblr}{colspec={Q[c,m,wd=9mm] X[l,m] Q[c,m,wd=10mm] Q[c,m,wd=9mm] Q[c,m,wd=9mm]}, width=\linewidth, rowhead=1, colsep=3pt,
  row{1}={bg=bkVol!14, font=\sffamily\bfseries\scriptsize}, hlines={0.4pt,bkLine}, rowsep=0.7pt}
章 & できるようになったこと & ページ & 1回め & 2回め \\
{\sffamily\bfseries\color{\bkTT}9-1} & $\sin$ の波をたして方形波に近づくのを表計算で確かめ、スペクトルを読める。 & \pageref{bk:9-1:spec} & \bkcb & \bkcb \\
 & フーリエ級数の形を書き、方形波・のこぎり波の係数を、ふつうの定積分で求められる。 & \pageref{bk:9-1:kakete} & \bkcb & \bkcb \\
 & 棒の熱と弦のゆれを、変数分離とフーリエ級数で解ける。 & \pageref{bk:9-1:netsu} & \bkcb & \bkcb \\
{\sffamily\bfseries\color{\bkTT}9-2} & ちがう周波数の $\sin$ どうしの積を1周期でたすと0になることを、表計算とグラフで確かめられる。 & \pageref{bk:9-2:kakeru} & \bkcb & \bkcb \\
 & 直交性を使って、係数が「かけて積分し、長さの2乗でわる」で取り出せるわけを説明できる。 & \pageref{bk:9-2:toridasu} & \bkcb & \bkcb \\
 & 場合分けの表9-2 と流れ図で、波の形に合った楽な求め方を選び、係数を計算できる。 & \pageref{bk:9-2:basho} & \bkcb & \bkcb \\
{\sffamily\bfseries\color{\bkTT}9-3} & 反対向きに回る2本の矢印をたすと $\cos$ の波になることを作図で確かめ、「負の周波数」の意味を言える。 & \pageref{bk:9-3:ya} & \bkcb & \bkcb \\
 & $a_n$・$b_n$ から複素フーリエ係数 $c_n=\frac{a_n-jb_n}{2}$ を求め、両側のスペクトルを読める。 & \pageref{bk:9-3:kaki} & \bkcb & \bkcb \\
 & 場合分けの表を使って $c_n$ を $e$ の定積分で求め、回路に入れた周期の波や、輪の形の針金の熱を計算できる。 & \pageref{bk:9-3:basho} & \bkcb & \bkcb \\
{\sffamily\bfseries\color{\bkTT}9-4} & 口笛・声・手をたたく音のスペクトルのちがいを説明でき、短いパルスほどスペクトルが広がることを言える。 & \pageref{bk:9-4:oto} & \bkcb & \bkcb \\
 & フーリエ変換 $F(\omega)$ の記号を読み、方形パルスの変換をふつうの定積分で求め、表計算の和で確かめられる。 & \pageref{bk:9-4:teigi} & \bkcb & \bkcb \\
 & 場合分けの表を使って変換を求め、DFT・FFT の結果を読み、サンプリング定理とエイリアシングを説明できる。 & \pageref{bk:9-4:basho} & \bkcb & \bkcb \\
{\sffamily\bfseries\color{\bkTT}9-5} & 「微分→$s$ をかける」の辞書を使って、冷めるお湯の微分方程式を $s$ の式に書きかえ、表でもどして答えを出せる。 & \pageref{bk:9-5:sekai} & \bkcb & \bkcb \\
 & ラプラス変換の記号を「名前」として読み、辞書（変換表）の行を、定義の積分で1回だけ作れる。 & \pageref{bk:9-5:gyou} & \bkcb & \bkcb \\
 & 微分の行 $sF(s)-f(0)$ を部分積分で導き、場合分けの表9-6（A・B・E）を使って微分方程式を変換できる。 & \pageref{bk:9-5:bibun} & \bkcb & \bkcb \\
{\sffamily\bfseries\color{\bkTT}9-6} & 表計算で $f(t)e^{-st}$ を台形でたし、変換表の値と同じになることを確かめられる。 & \pageref{bk:9-6:menseki} & \bkcb & \bkcb \\
 & $t^n$・$\sin$・$\cos$ と、$e^{-at}$ がかかった形の変換表の行を、部分積分・オイラーの公式・$s$ のずらしで作れる。 & \pageref{bk:9-6:tn} & \bkcb & \bkcb \\
 & 場合分けの表9-6と流れ図を使って、おくれ・パルス・デルタ関数・たたみこみをふくむ関数を変換し、初期値・最終値の定理で確かめられる。 & \pageref{bk:9-6:baai} & \bkcb & \bkcb \\
{\sffamily\bfseries\color{\bkTT}9-7} & 微分方程式を「変換する・わり算で解く・部分分数に分ける・表でもどす」の4つの手順で解く流れを、図とオイラー法の表で説明できる。 & \pageref{bk:9-7:tejun} & \bkcb & \bkcb \\
 & 分母がちがう1次式のかけ算になる式を、部分分数に分けて、表を逆に引いて時間の関数にもどせる。 & \pageref{bk:9-7:kurumaA} & \bkcb & \bkcb \\
 & 場合分けの表9-7を使って、重解・分母が0になる $s$ が複素数の式・入力のある式・おくれて始まる式・最後の値だけの場合も、もどし方を選んで計算できる。 & \pageref{bk:9-7:basho} & \bkcb & \bkcb \\
{\sffamily\bfseries\color{\bkTT}9-8} & 「入れると出てくる」箱を、伝達関数 $G(s)=Y(s)\div U(s)$ で表せることを、ブロック線図と表計算で説明できる。 & \pageref{bk:9-8:hako} & \bkcb & \bkcb \\
 & 微分方程式から伝達関数を作り、箱を直列・並列につないだ伝達関数と、ステップ応答・インパルス応答を求められる。 & \pageref{bk:9-8:tsukuru} & \bkcb & \bkcb \\
 & 場合分けの表9-8を使って、フィードバック・状態方程式・周波数応答（dB）の場合も、伝達関数を求めて使える。 & \pageref{bk:9-8:basho} & \bkcb & \bkcb \\
{\sffamily\bfseries\color{\bkTT}9-9} & 伝達関数の分母が0になる $s$（極）を×で打ち、その場所から動き方を当てられる。 & \pageref{bk:9-9:kyoku} & \bkcb & \bkcb \\
 & 1次おくれ・2次おくれの極から、時定数、行きすぎ、最初の山の時刻を求められる。 & \pageref{bk:9-9:niji} & \bkcb & \bkcb \\
 & すべての極の実部が負かどうかで安定を判定し、フィードバックで動く極、3次の判定、残るずれを計算できる。 & \pageref{bk:9-9:antei} & \bkcb & \bkcb \\
{\sffamily\bfseries\color{\bkTT}9-10} & 手のひらに立てた棒の「かたむき」と「たおれる速さ」を状態として、状態の点の動きの図を読める。 & \pageref{bk:9-10:jotai} & \bkcb & \bkcb \\
 & ばね・水そう・回路・倒立振子の式を状態方程式に書き、$A$ の固有値で落ち着くかを判定できる。 & \pageref{bk:9-10:tateru} & \bkcb & \bkcb \\
 & 状態フィードバックで固有値を選び、コンピュータの中のくり返しで落ち着く条件も判定できる。 & \pageref{bk:9-10:fb} & \bkcb & \bkcb \\
{\sffamily\bfseries\color{\bkTT}9-11} & オイラー法のくり返しを表計算で計算し、行列の指数関数 $e^{At}$ に近づくのを確かめられる。 & \pageref{bk:9-11:kiwami} & \bkcb & \bkcb \\
 & 対角の行列と対角化で $e^{At}$ を計算し、状態方程式を解ける。 & \pageref{bk:9-11:taikakuka} & \bkcb & \bkcb \\
 & 場合分けの表で $e^{At}$ の求め方を選び、固有値で落ち着くかを判定できる。 & \pageref{bk:9-11:basho} & \bkcb & \bkcb \\
{\sffamily\bfseries\color{\bkTT}9-12} & 音楽・写真の圧縮で、周波数の成分に分けて小さな成分をけずっていることを説明できる。 & \pageref{bk:9-12:asshuku} & \bkcb & \bkcb \\
 & 通信・MRI・機械の診断で、FFT と直交性が働いていることを言える。 & \pageref{bk:9-12:ofdm} & \bkcb & \bkcb \\
 & PID 制御の3つの項の役割を、極と最終値の定理のことばで説明できる。 & \pageref{bk:9-12:pid} & \bkcb & \bkcb \\
{\sffamily\bfseries\color{\bkTT}9-13} & 7本の路線の名前と終点を言い、全9巻の路線図でたどれる。 & \pageref{bk:9-13:rosen} & \bkcb & \bkcb \\
 & 路線が合流した5つの所を、式と例で説明できる。 & \pageref{bk:9-13:goryu} & \bkcb & \bkcb \\
 & 専門書の式を、名前の段で読み、記号を訳し、4つの道具で計算できる。 & \pageref{bk:9-13:tejun} & \bkcb & \bkcb \\
\end{longtblr}
\endgroup
